\documentclass[11pt,a4paper]{article}
\usepackage[utf8]{inputenc}
\usepackage[T1]{fontenc}
\usepackage[english]{babel}
\usepackage[margin=2.5cm]{geometry}
\usepackage{xcolor}
\usepackage{mdframed}
\usepackage{enumitem}
\usepackage{hyperref}
\hypersetup{colorlinks=true, urlcolor=blue, linkcolor=black}

% Gaya kotak komentar reviewer & jawaban penulis
\definecolor{revbg}{HTML}{EEF3F8}
\definecolor{ansac}{HTML}{2E7D32}
\newmdenv[backgroundcolor=revbg, linecolor=revbg, skipabove=6pt, skipbelow=4pt,
  innertopmargin=6pt, innerbottommargin=6pt]{revbox}
\newcommand{\reviewer}[1]{\vspace{4pt}\noindent\textbf{\textcolor{black}{#1}}}
\newcommand{\authors}[1]{\vspace{2pt}\noindent\textcolor{ansac}{\textbf{Author response:}} #1\par}
\newcommand{\loc}[1]{\vspace{1pt}\noindent{\small\textit{Location of change: #1}}\par\vspace{6pt}}

\title{\textbf{Response to Reviewers}\\[4pt]
\large Manuscript ID \#42551\\
\normalsize ``Design of a cloud backup system based on multi-compression algorithms
for ransomware attack mitigation''}
\author{Hero Yudo Martono, Nafisah Rahmadani Harahap\\
\small Politeknik Elektronika Negeri Surabaya, Indonesia}
\date{\today}

\begin{document}
\maketitle

\noindent We sincerely thank the Editor and all reviewers for their constructive and
detailed feedback. We have carefully revised the manuscript to address every comment.
Below we respond to each comment point by point. Reviewer comments are shown in shaded
boxes; our responses and the exact location of the corresponding changes follow each comment.

\medskip
\noindent\textbf{Summary of major revisions:}
\begin{itemize}[leftmargin=1.4em, itemsep=1pt]
  \item Clarified the dataset size discrepancy: the dataset consists of \textbf{43 files}
        consistently throughout the manuscript (the earlier value ``113'' has been corrected).
  \item The abstract now explicitly states that ransomware attacks were \textbf{simulated and
        controlled} (no real malware executed) on a \textbf{limited synthetic dataset}.
  \item Added a careful discussion explaining that the 0\% FPR/FNR is a logical consequence
        of hash-based integrity checking under controlled conditions, not a general claim.
  \item Corrected the technical order to \textbf{compress-then-encrypt} and re-ran all experiments.
  \item Added a State-of-the-Art paragraph, a research-gap paragraph with the entropy/compression
        formulation, and a comparison with detection-oriented literature.
  \item Validated the system on \textbf{two platforms} (Windows and Linux) on AWS EC2, with a
        cross-platform consistency table.
  \item Added a \textbf{statistical analysis} (5 repeated trials): mean~$\pm$~standard deviation of
        compression time, throughput, and CPU/RAM utilization.
  \item Expanded future work beyond Random Forest.
\end{itemize}

% =====================================================================
\section*{Response to the Editor}
% =====================================================================

\reviewer{Editor comment 1. Strengthen the structure, clarity, and argumentation; begin with a
strong thesis statement; each paragraph should have a topic sentence and logical flow.}
\begin{revbox}\small
Improve overall structure and argumentation; start the introduction with a clear thesis and
support it with solid evidence.
\end{revbox}
\authors{The Introduction has been restructured. It now opens with the escalating ransomware
threat (supported by quantitative statistics), builds a logical chain from the problem
(vulnerable connected backups, single-algorithm compression) to the research gap, and closes
with an explicit statement of contributions.}
\loc{Section~1 (Introduction), revised paragraphs and new State-of-the-Art paragraph.}

\reviewer{Editor comment 2. Highlight the state of the art in the Introduction; compare findings
with prior work in Results \& Discussion; define the contribution explicitly.}
\begin{revbox}\small
Emphasize state of the art; compare with previous results; state the novelty clearly.
\end{revbox}
\authors{We added a dedicated State-of-the-Art paragraph summarizing recent ransomware detection
and mitigation research (ensemble ML $>$99\%, deep learning 96--99\%, real-time detectors), and
we explicitly position our contribution at the \emph{mitigation/recovery} layer rather than
competing on detection accuracy. Results are now compared with the detection-oriented literature.}
\loc{Section~1 (State-of-the-Art paragraph) and Section~2 (Related Work, added comparison paragraph).}

\reviewer{Editor comment 3. The Method must include standard and novel approaches with
justification of validity, and reproducible step-by-step procedures.}
\begin{revbox}\small
Method should be reproducible and justify validity of the approach.
\end{revbox}
\authors{The Method now describes the corrected \emph{compress-then-encrypt} workflow with
justification (encrypted high-entropy data cannot be compressed). All steps (SHA-256 baseline,
compression, AES-256 encryption, air-gapped storage, detection, recovery) are described in order,
and the full implementation was executed on AWS EC2, making the procedure reproducible.}
\loc{Section~3 (Method), subsections on architecture, encryption, compression, and backup/recovery.}

\reviewer{Editor comment 4. The paper lacks critical discussion, comparison, and interpretation
of the implications.}
\begin{revbox}\small
Add critical discussion: what are the implications, and what is useful going forward?
\end{revbox}
\authors{We added critical discussion of the FPR/FNR results (their meaning and limitations), a
cross-platform consistency analysis, and a substantially expanded Conclusion covering limitations
and future directions.}
\loc{Section~4 (Results \& Discussion: FPR/FNR discussion, cross-platform subsection) and Section~5 (Conclusion).}

% =====================================================================
\section*{Response to Reviewer A}
% =====================================================================

\reviewer{A1. Strategic questions: scope fit, objective/approach, novelty, contribution to
knowledge, future implications, and why the journal should publish it.}
\begin{revbox}\small
How does the manuscript fit the journal scope and why is it of interest? What is the objective
and approach? What is the new contribution? How does it add to existing knowledge? What are the
future implications? Why should this journal publish it?
\end{revbox}
\authors{The manuscript fits the journal scope (secure and efficient cloud storage systems). Its
objective is an integrated, cost-effective ransomware-resilient cloud backup. The novelty is the
integration of logical air-gapped isolation, adaptive multi-algorithm compression, and
multi-indicator integrity detection in a single framework---an integration that prior studies
addressed only separately. These points are now stated explicitly in the Introduction and
Related Work.}
\loc{Section~1 (contribution paragraph) and Section~2 (novelty paragraph).}

\reviewer{A2. Abstract does not state that the ransomware attacks were simulated/controlled
(not real malware), nor that the dataset is synthetic and limited.}
\begin{revbox}\small
Abstract should state the attacks were simulated and controlled, and the dataset synthetic/limited.
\end{revbox}
\authors{The abstract now explicitly states that validation used a \emph{controlled simulated
attack model (no real malware executed)} on a \emph{limited synthetic dataset of 43 files}.}
\loc{Abstract (revised).}

\reviewer{A3. Introduction should emphasize that prior studies treated backup security,
ransomware mitigation, and compression separately, and did not analyze ransomware-resistant
backup architectures, immutable storage, snapshot recovery, deduplication, or zero-trust backup.}
\begin{revbox}\small
Clarify the research gap relative to immutable storage, snapshot recovery, deduplication,
zero-trust backup.
\end{revbox}
\authors{A research-gap paragraph was added, explicitly noting that prior work addressed these
aspects separately, and positioning our integrated framework against them. Related concepts
(immutable/WORM, zero-trust backup, deduplication) are now referenced and included in future work.}
\loc{Section~1 (research-gap paragraph) and Section~5 (future work).}

\reviewer{A4. Method issues: (i) the workflow is confusing; (ii) adaptive compression is only
rule-based; (iii) CVE mapping lacks technical rigor; (iv) detection relies too heavily on hash
and metadata changes.}
\begin{revbox}\small
Clarify workflow; state rule-based compression explicitly; strengthen CVE mapping; discuss the
reliance on hash/metadata.
\end{revbox}
\authors{(i) The workflow was rewritten in clear sequential steps with the corrected
compress-then-encrypt order. (ii) The compression selection is now explicitly described as
\emph{rule-based} (by file type and size). (iii) Each simulated attack is mapped to a CVE pattern
(CVE-2017-0144 encrypt-and-hold, CVE-2018-8453 metadata/header corruption, CVE-2021-36934
overwrite-and-corrupt) with a description of the controlled file modification performed. (iv) We
now explicitly acknowledge that hash/metadata-based detection is an integrity mechanism (detects
\emph{that} a file changed) and discuss its limitations against advanced evasion (FPE, fileless, LotL).}
\loc{Section~3 (Method: architecture, compression, attack simulation, detection) and Section~4 (FPR/FNR discussion).}

\reviewer{A5. Results lack detailed statistical analysis (repeated trials, standard deviation,
CPU/RAM utilization, throughput, energy, cloud cost) and comparison with deduplication/incremental
backup; some figures are not referenced.}
\begin{revbox}\small
Add statistical analysis and comparisons; ensure all figures are referenced.
\end{revbox}
\authors{We have substantially strengthened the statistical analysis. We now report, in a new
subsection, \textbf{repeated-trial statistics (5 trials)} for every algorithm on the full dataset:
mean and standard deviation of compression time, throughput (MB/s), and CPU/RAM utilization
(measured with \texttt{psutil}) on AWS EC2. The very small standard deviations (e.g., Snappy
$0.003$~s, Zstandard $0.016$~s) confirm stable performance, and the throughput ranges (Snappy/LZ4
$>590$~MB/s vs.\ Gzip $30.5$~MB/s) provide an objective basis for the adaptive selection. All figures
are now referenced in the text before they appear. In addition, we added a \textbf{quantitative
comparison of the proposed adaptive strategy against baselines} (no-compression, static-Gzip,
static-Zstandard) on the full dataset. This reveals a clear speed-vs-storage trade-off: the adaptive
strategy is about $21\times$ faster than static Gzip ($0.87$~s vs.\ $18.84$~s; $611.8$ vs.\
$28.2$~MB/s) while still achieving $43.3\%$ storage saving, showing that its main advantage is the
optimal balance between storage efficiency and processing speed for low-RTO backup. Energy profiling
and direct comparison with deduplication/incremental backup remain listed as future work.}
\loc{Section~4 (statistical-analysis subsection and new baseline-comparison subsection with tables; all figure references) and Section~5 (remaining future work).}

\reviewer{A6. Conclusion future work is too narrow (only Random Forest).}
\begin{revbox}\small
Broaden future work beyond Random Forest.
\end{revbox}
\authors{Future work now includes: ML-based compression selection, testing with real ransomware
in a safe sandbox, immutable/WORM and zero-trust backup integration, formal security analysis and
key management, enterprise-scale evaluation (throughput, CPU/RAM, energy, cloud cost), and direct
comparison with incremental/deduplication backup.}
\loc{Section~5 (Conclusion, expanded future work).}

\reviewer{A7. Ensure references [21], [41]--[44] are cited in the text; include DOIs where available.}
\begin{revbox}\small
Cite all references in the text and add DOIs.
\end{revbox}
\authors{We reviewed all references to ensure each is cited in the text. DOIs are being added for
all sources that have them (e.g., the Kapoor et al. survey, \texttt{doi:10.3390/su14010008}). Any
reference that turns out to lack a verifiable DOI or complete bibliographic data will be replaced
with a properly indexed source.}
\loc{References section (updated entries and in-text citations).}

% =====================================================================
\section*{Response to Reviewer B}
% =====================================================================

\reviewer{B1. The abstract should state that the system was validated using a simulated attack
model on a synthetic dataset of 43 files.}
\begin{revbox}\small
State validation on a simulated model with a synthetic dataset of 43 files in the abstract.
\end{revbox}
\authors{Done. The abstract now states the controlled simulation and the 43-file synthetic dataset explicitly.}
\loc{Abstract (revised).}

\reviewer{B2. The Introduction needs a paragraph on the mathematical problem of compression
efficiency for high-entropy (encrypted) vs low-entropy (plaintext) data, and a defined research gap.}
\begin{revbox}\small
Add a paragraph on the entropy/compression mathematical problem and define the research gap.
\end{revbox}
\authors{We added a paragraph based on Shannon entropy, $H(X)=-\sum_i p(x_i)\log_2 p(x_i)$,
explaining the compression limit and why high-entropy (encrypted/already-compressed) data barely
compresses, while low-entropy text compresses well. This directly motivated the corrected
compress-then-encrypt order and the defined research gap.}
\loc{Section~1 (research-gap/entropy paragraph).}

\reviewer{B3. Figures 1, 2, 3, 10 must be referenced and explained in the text before they appear;
Figures 14, 15, 16 were not found.}
\begin{revbox}\small
Reference and explain all figures before display; restore missing Figures 14--16.
\end{revbox}
\authors{All figures are now referenced and explained in the text prior to display. The full set
of figures (including the recovery/restore figures) has been regenerated from the actual
experimental outputs and is included in the revised manuscript.}
\loc{Sections~3 and~4 (all figure references and captions).}

\reviewer{B4. Tables 3, 9, 11 must be referenced and explained in the text before they appear.}
\begin{revbox}\small
Reference and explain Tables 3, 9, 11 before display.
\end{revbox}
\authors{All tables are now introduced and explained in the text before they appear, including the
dataset table (with an explicit total of 43 files), the hash/metadata-change tables, and the
detection-accuracy table.}
\loc{Section~3 (dataset table) and Section~4 (result tables).}

\reviewer{B5. The Conclusion should be more concise: summarize key findings, answer the objectives,
and state contribution and future work.}
\begin{revbox}\small
Make the conclusion more concise and goal-oriented.
\end{revbox}
\authors{The Conclusion was rewritten to summarize the key quantitative findings, confirm the
objectives, state the contribution, and clearly list future work.}
\loc{Section~5 (Conclusion).}

\reviewer{B6. Reference [21] is not cited in the text; all references must be cited.}
\begin{revbox}\small
Ensure reference [21] and all others are cited.
\end{revbox}
\authors{We verified that all references, including the NIST storage-security guideline, are now
cited in the text.}
\loc{References and in-text citations.}

% =====================================================================
\section*{Response to Reviewer C (score: 8)}
% =====================================================================

\reviewer{C1. English grammar contains errors (missing articles/subjects/auxiliary verbs), e.g.,
``evaluate the resilience\ldots'' should be ``To evaluate the resilience\ldots''.}
\begin{revbox}\small
Perform a thorough English language edit.
\end{revbox}
\authors{The manuscript underwent a thorough language edit to correct grammar, add missing
articles/subjects, and improve consistency and readability throughout.}
\loc{Throughout the manuscript.}

\reviewer{C2. Dataset size ambiguity: the description gives 9 CSV + 9 LOG/TXT + 14 JPG + 11 XLSX
= 43 files, but the text repeatedly reports restoring 113 files.}
\begin{revbox}\small
Clarify the 43 vs 113 file discrepancy.
\end{revbox}
\authors{This has been fully corrected. The dataset is \textbf{43 files} consistently. Every prior
occurrence of ``113'' was corrected to 43 (abstract, narrative, RTO table, and detection-accuracy
table), and a Total row (43) was added to the dataset table. The re-run experiments confirm
recovery of 43/43 files.}
\loc{Abstract, Section~3 (dataset table), Section~4 (RTO and detection tables).}

\reviewer{C3. The 0\% FPR and 0\% FNR results are overstated; such perfection is rare in
real-world ransomware detection.}
\begin{revbox}\small
Discuss the 0\% FPR/FNR result more cautiously.
\end{revbox}
\authors{We added an explicit discussion clarifying that 0\% FPR/FNR is a \emph{logical consequence}
of deterministic hash-based integrity checking under controlled conditions (any 1-byte change
alters the SHA-256 hash), and \textbf{not} a claim of real-world detection perfection. We state the
limitations against advanced evasion (FPE, fileless, LotL) and frame the metric as functional
validation of the mitigation/recovery layer.}
\loc{Section~4 (FPR/FNR discussion paragraph).}

\reviewer{C4. Describe the compression selection mechanism explicitly as rule-based.}
\begin{revbox}\small
State clearly that compression selection is rule-based.
\end{revbox}
\authors{The compression selection is now explicitly described as a \emph{rule-based} strategy
driven by file type and size, and this is stated in the Introduction, Method, and Conclusion.}
\loc{Sections~1, 3, and~5.}

\reviewer{C5. Strengthen experimental validation with additional statistical analysis.}
\begin{revbox}\small
Add statistical analysis.
\end{revbox}
\authors{We strengthened the empirical basis in two ways. First, we re-ran the full experiment on
two platforms (Windows Server 2022 and Ubuntu 22.04 on AWS EC2) and added a cross-platform
consistency table. Second, we added a dedicated statistical analysis with \textbf{5 repeated trials}
per algorithm, reporting mean~$\pm$~standard deviation of compression time, throughput, and CPU/RAM
utilization (Table on compression statistics). These measurements confirm the stability and
reproducibility of the results. For clarity and honest reporting, we now state explicitly that the
local Windows machine is the \emph{initial development environment}, while the AWS EC2 environment
is the \emph{primary measurement/validation environment} from which all reported quantitative
results are obtained; the original Windows setup is retained (not removed) to preserve continuity
with the initially submitted version.}
\loc{Section~4 (cross-platform subsection, consistency table, and new statistical-analysis subsection).}

% =====================================================================
\section*{Additional journal requirements}
% =====================================================================
\authors{We have also completed the Funding, Author Contributions (CRediT), Conflict of Interest,
and Data Availability sections, and included author biographies. The revision is submitted under
the same manuscript ID.}

\vspace{6pt}
\noindent We again thank the Editor and reviewers for their valuable feedback, which has
substantially improved the manuscript.

\end{document}
